\NativeTitle{25}{第四单元素养评价（A）}{（运算律）}{45}{4}{一 & 二 & 三 & 四}{ &  &  & }
\NativeSlot{283}{36}{%
\ExamGroup{一、填空。{\normalfont\kaishu\fontsize{10.8bp}{14bp}\selectfont （第4题6分，其余每空1分，共21分）}}

% p025_s01

\ExamQuestion{1.}{5}%
\(a\times(b+c)=\)\mbox{（\hspace{1.8em}）}\(\times\)\mbox{（\hspace{1.8em}）}\(+\)\mbox{（\hspace{1.8em}）}\(\times\)\mbox{（\hspace{1.8em}）}。\par
}
\NativeSlot{332}{56}{%
% p025_s02

\ExamQuestion{2.}{5}%
观察右边的竖式，在计算过程中运用了\mbox{（\hspace{1.8em}）}律，用横式表示右边竖式的计算过程是\mbox{（\hspace{1.8em}）}。\par
\BookFigure{figures/p025_s02-01.tex}

% p025_s03

\ExamQuestion{3.}{5}%
小明在用计算器计算\(87\times99\)时，发现按键“9”出现故障，聪明的你想怎样用计算器算出结果？请写出你的算式：\mbox{（\hspace{1.8em}）}。\par
}
\NativeSlot{434}{68}{%
\ExamQuestion{4.}{5}%
在方框里填上合适的数或字母，在圆圈里填上“\(+\)、\(-\)、\(\times\)、\(\div\)”。\par\smallskip
\noindent\ExamTile{.49}{%
% p025_s04
\(98+74+26=\square+(74+\square)\)\par
% p025_s05
\(m\times21\times9=m\times(\square\bigcirc\square)\)\par
% p025_s06
\(380\div5\div2=380\bigcirc(\square\bigcirc\square)\)\par
}\hfill\ExamTile{.49}{%
% p025_s07
\(612-124-76=612\bigcirc(\square\bigcirc\square)\)\par
% p025_s08
\(389+102=389\bigcirc\square\bigcirc\square\)\par
% p025_s09
\((a+b)\times9=a\times\square+\square\times\square\)\par
}\par
}
\NativeSlot{553}{71}{%
\ExamQuestion{5.}{5}%
在圆圈里填上“\(>\)”“\(<\)”或“\(=\)”。\par\smallskip
\noindent%
% p025_s10
\makebox[.499\linewidth][l]{\(50\times(12+33)\ \bigcirc\ 50\times12+33\)}%
% p025_s11
\makebox[.499\linewidth][l]{\(132-68-12\ \bigcirc\ 132-68+12\)}\par
\noindent%
% p025_s12
\makebox[.499\linewidth][l]{\(720\div8\times6\ \bigcirc\ 720\div8\div6\)}%
% p025_s13
\makebox[.499\linewidth][l]{\(32\times20\times8\ \bigcirc\ 32\times(20\times8)\)}\par
\noindent%
% p025_s14
\makebox[.499\linewidth][l]{\(12\times56\ \bigcirc\ 56\times12\)}%
% p025_s15
\makebox[.499\linewidth][l]{\(85\times103\ \bigcirc\ 85\times100+3\)}\par
}
\NativeSlot{672}{36}{%
% p025_s16

\ExamQuestion{6.}{5}%
如果\(\bigcirc+\triangle=4\)，那么\(25\times\bigcirc+25\times\triangle=\)\mbox{（\hspace{1.8em}）}。\par
}
\NativeSlot{701}{36}{%
% p025_s17

\ExamQuestion{7.}{5}%
明明在计算时，把\(8\times(\square+6)\)错算成\(8\times\square+6\)，他得到的结果与正确结果相差\mbox{（\hspace{1.8em}）}。\par
}
